\section*{Capítulo \ref{ch:b3:group-theory-applied} --- Teoria de grupos aplicada}

\begin{solution}{exo:b3:group-theory-applied:1}
$n_{A_1} = \frac14(5 - 1 + 3 + 1) = 2$, $n_{A_2} = \frac14(5 - 1 - 3 - 1) = 0$,
$n_{B_1} = \frac14(5 + 1 + 3 - 1) = 2$, $n_{B_2} = \frac14(5 + 1 - 3 + 1) = 1$:
$\Gamma = 2A_1 + 2B_1 + B_2$, de dimensão 5.
\end{solution}

\begin{solution}{exo:b3:group-theory-applied:2}
Mesma contagem de átomos deixados no lugar que a água: $\Gamma_{\mathrm{vib}} = 2A_1 + B_1$. $A_1$ ($z$) e
$B_1$ ($x$) são \omterm{def:b3:group-theory-applied:activity}{ativos no IV}, e ambos são também \omterm{def:b3:group-theory-applied:activity}{ativos no Raman} ($x^2$, $xz$): três
bandas em cada espectro.
\end{solution}

\begin{solution}{exo:b3:group-theory-applied:3}
$B_1\otimes B_2$: \omterm{def:b3:point-groups:representation}{caracteres} $(1, 1, -1, -1)$ $= A_2$. $A_2\otimes B_1$: $(1, -1, -1,
1) = B_2$. Em $C_{3v}$, $E\otimes E$ tem \omterm{def:b3:point-groups:representation}{caracteres} $(4, 1, 0)$; redução: $A_1 + A_2 +
E$.
\end{solution}

\begin{solution}{exo:b3:group-theory-applied:4}
Infravermelho: os dois modos $t_2$, 3019 e \qty{1306}{cm^{-1}}. Raman: os quatro ($a_1$,
$e$, $t_2$ se transformam todos como funções quadráticas). Sem \omterm{def:b3:point-groups:inversion}{centro de inversão},
as coincidências são permitidas.
\end{solution}

\begin{solution}{exo:b3:group-theory-applied:5}
Átomos no lugar $(4, 1, 2, 4, 1, 2)$ vezes $\chi_{xyz} = (3, 0, -1, 1, -2, 1)$ dão
$\chi_{3N} = (12, 0, -2, 4, -2, 2)$, que se reduz a $A_1' + A_2' + 3E' + 2A_2'' +
E''$. Retire $E' + A_2''$ (translações) e $A_2' + E''$ (rotações): $A_1' + 2E'
+ A_2''$. IV: $2E' + A_2''$, três bandas; Raman: $A_1' + 2E'$, três bandas. $A_2''$
(deformação fora do plano) aparece só no IV, $A_1'$ (estiramento simétrico) só no
Raman, e os modos $E'$ aparecem nos dois.
\end{solution}

\begin{solution}{exo:b3:group-theory-applied:6}
IV: os dois modos $T_{1u}$, duas bandas. Raman: $A_{1g}$, $E_g$, $T_{2g}$, três bandas.
$T_{2u}$ é silencioso. \ce{SF6} tem \omterm{def:b3:point-groups:inversion}{centro de inversão}: exclusão mútua.
\end{solution}

\begin{solution}{exo:b3:group-theory-applied:7}
Sob $E$, $C_3$, $C_3^2$, a função $h_2$ vai a $h_2$, $h_3$, $h_1$ (\omterm{def:b3:point-groups:representation}{caracteres} 2,
$-1$, $-1$); as reflexões contribuem com 0. $\hat P_Eh_2 \propto 2h_2 - h_3 - h_1$.
Sua sobreposição com $2h_1 - h_2 - h_3$ (desprezando as sobreposições atômicas) é $-3$, e
esta última tem norma $\sqrt6$; subtraindo a projeção, $(2h_2 - h_1 - h_3) +
\frac12(2h_1 - h_2 - h_3) = \frac32(h_2 - h_3)$.
\end{solution}

\begin{solution}{exo:b3:group-theory-applied:8}
fac ($C_{3v}$, \omterm{def:b3:point-groups:representation}{caracteres} $3, 0, 1$): $A_1 + E$, duas bandas no IV. mer ($C_{2v}$,
\omterm{def:b3:point-groups:representation}{caracteres} $3, 1, 3, 1$): $2A_1 + B_1$, três bandas. \ce{M(CO)6} ($O_h$): $A_{1g} + E_g
+ T_{1u}$, uma banda no IV ($T_{1u}$).
\end{solution}

\begin{solution}{exo:b3:group-theory-applied:9}
\omterm{def:b3:point-groups:representation}{Caracteres} $6, 0, 0, 2, 2, 0, 0, 0, 4, 2$: só $E$, os $C_4$ e $C_2$ ao longo dos
eixos (dois ligantes deixados no lugar), os três $\sigma_h$ (quatro) e os seis $\sigma_d$
(dois) deixam ligantes no lugar. A redução dá $A_{1g} + E_g + T_{1u}$. Os orbitais $s$
($a_{1g}$), $d_{z^2}$ e $d_{x^2-y^2}$ ($e_g$) e $p_x$, $p_y$, $p_z$ ($t_{1u}$) do metal se
combinam; $t_{2g}$ permanece não ligante.
\end{solution}

\begin{solution}{exo:b3:group-theory-applied:10}
$n_i = \frac1h(1\cdot h\cdot\chi_i(E)) = d_i$. A dimensão da representação
regular é $h$, e também $\sum_in_id_i = \sum_id_i^2$.
\end{solution}

\begin{solution}{exo:b3:group-theory-applied:11}
Em $D_{2h}$ (molécula ao longo de $z$), átomos no lugar $(4, 4, 0, 0, 0, 0, 4, 4)$, $\chi_{3N}
= (12, -4, 0, 0, 0, 0, 4, 4)$, que se reduz a $2A_g + 2B_{2g} + 2B_{3g} + 2B_{1u} + 2B_{2u}
+ 2B_{3u}$. Retire as translações $B_{1u} + B_{2u} + B_{3u}$ e as duas rotações
$B_{2g} + B_{3g}$: $2A_g + B_{2g} + B_{3g} + B_{1u} + B_{2u} + B_{3u}$, sete modos.
Em $D_{\infty h}$: $2\Sigma_g^+ + \Pi_g + \Sigma_u^+ + \Pi_u$. IV: $\Sigma_u^+$ e $\Pi_u$
(duas bandas); Raman: $2\Sigma_g^+$ e $\Pi_g$ (três bandas); nenhuma coincidência.
\end{solution}

\begin{solution}{exo:b3:group-theory-applied:12}
$1b_2 \to 4a_1$: $B_2$, $y$, permitida, perpendicular ao plano. $3a_1 \to 4a_1$:
$A_1$, $z$, permitida, ao longo do eixo $C_2$. $1b_2 \to 2b_1$: $A_2$, proibida.
$1b_1 \to 4a_1$: $B_1$, $x$, permitida, no plano, perpendicular ao eixo.
\end{solution}

\begin{solution}{pb:b3:group-theory-applied:1}
\textbf{1.} cis: os dois CO em vértices adjacentes; trans: opostos; fac: três
CO em uma face; mer: três CO em um meridiano (dois em trans, um entre eles).
\textbf{2.} $C_{2v}$.
\textbf{3.} $D_{4h}$.
\textbf{4.} $C_{3v}$.
\textbf{5.} $C_{2v}$.
\textbf{6.} Apenas trans-\ce{ML4(CO)2}.
\textbf{7.} Uma operação que leva um CO sobre outro põe seu vetor fora da
diagonal (contribuição 0); uma que deixa um CO no lugar leva seu vetor de
ligação sobre ele mesmo (+1); nenhuma inverte um vetor C--O.
\textbf{8.} \omterm{def:b3:point-groups:representation}{Caracteres} $(2, 0, 2, 0)$, com $\sigma_v(xz)$ o plano dos dois CO: $A_1 +
B_1$.
\textbf{9.} \omterm{def:b3:point-groups:representation}{Caracteres} 2 sob $E$, $C_4$, $C_2$, $\sigma_v$, $\sigma_d$, 0 nas demais:
$A_{1g} + A_{2u}$.
\textbf{10.} $(3, 0, 1)$: $A_1 + E$.
\textbf{11.} $(3, 1, 3, 1)$: $2A_1 + B_1$.
\textbf{12.} 2, 2, 3, 3: uma dimensão por CO.
\textbf{13.} cis 2, trans 1, fac 2, mer 3 (estiramentos \omterm{def:b3:group-theory-applied:activity}{ativos no IV}).
\textbf{14.} Os mesmos números: cis 2, trans 1, fac 2, mer 3.
\textbf{15.} trans: sua banda no IV ($A_{2u}$) e sua banda Raman ($A_{1g}$) são
modos diferentes.
\textbf{16.} $A_{2u}$ troca de sinal sob $i$: um CO se alonga enquanto o outro
se encurta, fora de fase.
\textbf{17.} $E$ é um par degenerado: dois modos de mesma frequência dão uma
só banda, mais a banda $A_1$.
\textbf{18.} Um isômero trans-dissubstituído, centrossimétrico.
\textbf{19.} Três bandas no IV: mer (fac daria duas).
\textbf{20.} Sim: mer tem três modos CO \omterm{def:b3:group-theory-applied:activity}{ativos no Raman} nas mesmas frequências,
por não ser centrossimétrico; fac mostraria duas.
\textbf{21.} $A_1$: o estiramento em fase é totalmente simétrico.
\textbf{22.} Uma mistura fac/mer mostraria até cinco bandas no IV, com as bandas
do fac em outras posições; três bandas limpas correspondem a um único isômero.
Um espectro de RMN de ${}^{13}$C decidiria: dois sinais de carbonila na razão 2:1
para mer, um único sinal para fac.
\textbf{23.} \textbf{O isômero mer tem três estiramentos CO \omterm{def:b3:group-theory-applied:activity}{ativos no IV}, o isômero
fac, dois}: o produto é mer.
\end{solution}
